\begin{figure}[!t]
\centering
\begin{tikzpicture}[
  box/.style={draw,rounded corners=2pt,align=center,font=\scriptsize,
              minimum height=8mm,inner sep=2pt},
  arrow/.style={-{Latex[length=1.8mm]},thin},
  note/.style={font=\tiny,align=center,text=black!70}
]
\node[box,fill=blue!8,text width=1.45cm] (p1) {Evidence\\payload};
\node[box,fill=blue!14,text width=1.65cm,right=2.5mm of p1] (h256)
  {SHA-256\\32-byte digest};
\node[box,fill=green!12,text width=1.85cm,right=2.5mm of h256] (a256)
  {EVM anchor\\algorithm ID + digest};
\draw[arrow] (p1) -- (h256);
\draw[arrow] (h256) -- (a256);

\node[box,fill=orange!8,text width=1.45cm,below=8mm of p1] (p2)
  {Same evidence\\payload};
\node[box,fill=orange!16,text width=1.65cm,right=2.5mm of p2] (h512)
  {SHA-512\\64-byte digest};
\node[box,fill=green!12,text width=1.85cm,right=2.5mm of h512] (a512)
  {Same EVM anchor\\algorithm ID + digest};
\draw[arrow] (p2) -- (h512);
\draw[arrow] (h512) -- (a512);

\node[note,below=2mm of h512,text width=4.2cm] (boundary)
  {SHA-256 also has a standard EVM precompile.\\
   SHA-512 has no standard EVM precompile; its full digest is anchored.};
\end{tikzpicture}
\caption{Fair SHA-2 comparison boundary. Both algorithms hash identical bytes
off chain and anchor their complete digest on identical EVM bytecode. The
resulting network comparison measures digest anchoring, not encryption or CPU
hash throughput.}
\label{fig:sha2-paths}
\end{figure}
